% ------------------------------------------------------------------------
% file `main-bookexercise-2-solution-body.tex'
%
%     solution of type `bookexercise' with id `2'
%
% generated by the `exerciseanswer' environment of the
%   `xsim' package v0.21 (2022/02/12)
% from source `main' on 2026/10/07 on line 1278
% ------------------------------------------------------------------------
    1. 选ABD。\\

    2. 选C。栈的插入删除只能在栈顶进行，C错误。\\

    3. 选D。对于C项，队列仅要求队首元素可以出队，但没要求仅队首元素可以被读取。你当然可以为队列实现获取第i个元素的函数接口，但通常来说我们使用队列的场景往往不会有这个需求。
    额外提示：你当然可以为你的“队列”实现第i个元素“出队”的函数接口，但这样做的话，你所实现的数据结构就不能叫做\term{队列}了。\\

    4. 选D。2134可由入入出出入出入出达成；2341可由入入出入出入出出达成；3241可由入入入出出入出出达成；栈满足后进先出，如论如何4一定比3先出，2一定比1先出，D无法达成。\\

    5. 选B。对于A，顺序栈的Push和Pop操作均在顺序表的表尾实现，时间均可达到O(1)；对于C，仅有head指针的单链表实现“头删”，即出队，可达到O(1)，但要实现“尾删”，即出队，
    则需要先找到最后一个节点，时间为O(n)。单链表中额外记录一个尾指针rear，则入队和出队均可达成O(1)时间。对于D，循环队列的入队和出队无需移动元素。\\

    6. 选D。通过本题给了我们一个提示，对于循环队列来说，初始front和rear不一定非要指向0，因为从环的几何结构来讲，“队首元素”保存在0位置和保存在i位置是没有区别的。
    仅有初始状态（队空）会满足front == rear。\\

    7. 选A、C、D、G。对于B，无论栈中有多少个元素，一次出栈操作只需要访问当前栈顶元素，然后修改栈顶指针即可，时间复杂度为O(1)。\\
    对于C，可由两栈的栈顶下标直接分别计算得到两栈的元素数量，时间均为O(1)，再相加即可得到总数量，总体时间为O(1)，C正确；\\
    对于D，要注意：当栈S2中没有元素，而栈S1将空间全部占满时，同样可以用S1.top+1==S2.top判断空间满，当S2将空间占满时同样如此，D正确；\\
    E、F显然错误；\\
    对于G，当栈满时，S1.top指向S1栈顶元素的下一个位置，即S2的栈顶元素位置；对应地，S2.top指向S1的栈顶元素位置，S1.top>S2.top成立，G正确。\\

    8. 选C。队列的特性决定了第i个入队的元素一定是第i个出队的，而栈则不同，栈的出栈顺序是会受到操作过程影响的，{\color{red}\textbf{参考本习题第4题}}。\\

    9. 选AC。\\

    10. 选BCD。\\

    11. 选A。对于循环双链表，可以在O(1)时间立即拿到尾节点指针，因此在头、尾操作均为O(1)。\\

    12. 选C。A项可以通过\Circled{1}\Circled{1}\Circled{2}\Circled{2}\Circled{1}\Circled{1}\Circled{3}\Circled{3}得到；
    B项可通过\Circled{2}\Circled{1}\Circled{1}\Circled{1}\Circled{1}\Circled{1}\Circled{3}得到；
    D项输入与输出刚好相反，显然通过栈就能逆序：\Circled{2}\Circled{2}\Circled{2}\Circled{2}\Circled{2}\Circled{2}\Circled{3}\Circled{3}\Circled{3}\Circled{3}\Circled{3}\Circled{3}。
    对于C，3在1、2前输出，因此1、2必按顺序先入栈，所以1必晚于2输出，C错误。\\

    13. 选C。这是一个非常有意思的题。不妨假设仅有左边能出队。无论如何a一定先入队，摆在最中央的位置。接下来，每一个元素入队，你都可以选择将其“摆在队中已有元素的左边或右边”。所以你需要
    以a为中心努力地去构造四个选项所给出的序列。如A选项要求的序列是b,a,c,d,e，于是你将b摆在a左边，将c摆在ba的右边，将d摆在bac的右边，再将e摆在bacd的右边，便得到
    了A项所要求的序列，从左边出队，便得到了b,a,c,d,e。B、D项也都可以通过同样的方法得到。\\

    14. 选D。\\

    15. 选B。详细过程已经在\ref{sec:表达式的转换与求值}中介绍过，不再赘述。\\

    16. 选B。本题比较简单，注意出栈顺序别搞错，同时注意二元运算符op的操作数顺序是b op a而不是a op b，不要搞错就行。

